\section{Problem and Motivation}
\label{sec:problem}
\textbf{Problem.} A candidate $W=(w_1,w_2,w_3)$ holds one weight vector $w_b\in\Delta^6$ per building (non-negative,
summing to one). The simulator returns
\begin{equation}
R(W)=\frac{1}{8}\sum_{k=1}^{8}\sum_{t=2}^{168}\bigl|E_{k,t}(W)-E_{k,t-1}(W)\bigr|\ \text{kWh},
\label{eq:R}
\end{equation}
where $E_{k,t}$ is the district net load in hour $t$ of week $k$. The best training row (\texttt{p01969}) has
$R=115.76$, which is also the offline bar; the course quadratic-ridge baseline's optimum measures 119.09.

\textbf{Surrogates fail away from the data.} On a frozen ten-split protocol, our surrogates (quadratic ridge,
ExtraTrees, gradient boosting, a stacked ensemble) scored within one standard deviation of each other. Their optima,
however, lay far from the data, and the two points we could measure there contradicted them (Table~\ref{tab:offline}):
the baseline optimum was predicted at 75.9 and measured 119.09, and our first submission, the consensus of three tree
models, was predicted at 107--111 and measured 119.71. Our second submission was therefore a step of length 0.05 from
the best row, chosen by an additive ARD Gaussian process~\cite{rasmussen} as its most confident direction among
19,483 random ones. It measured 115.68, 0.08~kWh below the bar, but its gain ($-0.08$) was 35 times smaller than predicted.
The same holds online: the additive GP predicts the 200 lowest of 6,548 locally measured points ($R\approx105.1$) at
$124.4\pm8.8$, because its prior mean reverts to the data average beyond about two training spacings.

\begin{table}[t]
\centering
\caption{Surrogate predictions at the two points measured before our second offline submission (kWh).}
\label{tab:offline}
\begin{tabular}{lrr}
\toprule
Model & baseline optimum & submission 1\\
\midrule
quadratic ridge (course) & 75.90 & 95.13\\
HistGradientBoosting & 117.62 & 107.80\\
XGBoost & 114.21 & 107.28\\
stacked ensemble & 118.73 & 111.42\\
\midrule
\textbf{measured} & \textbf{119.09} & \textbf{119.71}\\
\bottomrule
\end{tabular}
\end{table}

\textbf{Structure that measurements reveal.} After the offline phase the simulator was released, and the TA allowed
local runs to design the policy, provided the submitted code uses only the training rows before its first call. An
exploratory search (about 22,500 local evaluations, never an input to the policy) showed three regularities. First,
$R$ is additive over buildings within a basin, i.e.\ the low-$R$ region around one combination of dominant objectives
(interaction error about 0.03~kWh). The training rows agree globally, more loosely: a building-additive model
cross-validates as well as an unconstrained one ($R^2$ 0.846 vs.\ 0.844), so additivity costs almost no fit. Second, $R$ is
terraced: along one weight it is constant on plateaus 0.02--0.2\,\% of a weight wide, so many moves return exactly
the same value. Third, tiny weights act as switches: shares of $10^{-5}$ to $3\times10^{-3}$ on otherwise unused
objectives flip near-tied hourly decisions, and the last gains came from one-coordinate scans. Additivity motivates
Section~\ref{sec:block}; terraces and the uneven payoff of move types motivate Section~\ref{sec:alloc}.
